\subsection{Unresolved feedback effects and event detection}\label{sec:live-study}
The separate pilot holds initial messages fixed and varies whether supplementary tool records enter the next model request. \emph{Full} forwards the whole response; \emph{Projected} retains task-required fields. The synthetic records deliberately carry action information. Two fixed 8B models perform 3,120 local LangGraph/MCP executions over streaming TLS~1.3/HTTP/2; each test condition contains 120 tasks in twenty nonempty and four empty sessions.

First-feedback LR recognition changes from 43 to 23 correct for Qwen and from 40 to 22 for Llama. All eight primary local comparisons, including post-feedback versus prefix, exceed .05 after Holm correction (minimum .0508). Complete-execution LR changes from 46 to 44 and from 62 to 33. Projection changes content, length and execution order together, so the pilot leaves the contribution of feedback content and end-to-end protection unresolved. The hosted-alias retest also has no corrected primary rejection. Supplement S1 retains every attack, phase, interval and execution outcome.

Concurrent runs additionally test continuous detectors with visible connection identities and no task windows. Their Shared-condition precision is 4.3--18.6\%, with 8.21--44.79 false events per test session at IoU .5. A constant all-active baseline ranks above learned detectors at IoU .3 and below them at .5. This diagnoses detector and matching-rule behavior in the constructed workload. Full .3/.5/.7 results, background definitions and mixed shaping outcomes remain in Supplement S1. These diagnostics guide evaluation design; they do not measure general Agent localization difficulty.
